\documentclass[11pt,a4paper]{article}
\usepackage[margin=1in]{geometry}
\usepackage{amsmath,amssymb,amsthm}
\usepackage[T1]{fontenc}\usepackage{lmodern}\usepackage{booktabs}
\usepackage[colorlinks=true,linkcolor=black,urlcolor=blue]{hyperref}
\newcommand{\Psip}{\Psi^{+}}\newcommand{\Ts}{\mathcal{T}}\newcommand{\Tp}{\mathcal{T}^{+}}
\title{\textbf{You are right about the two entries --- and your boxed identity
has a spurious sign}\\[2mm]
\large The correct relation is $\Psip(f)(u)=+\Psi(f\circ\varphi)(-u)$, with no sign
and no $n$-dependence. Your conclusion is unaffected.}
\author{Clio}\date{2026-09-01}
\begin{document}\maketitle
\begin{center}\begin{tabular}{@{}ll@{}}\toprule
\textbf{Author} & Clio\\ \textbf{Recipient} & Rick\\ \textbf{Date} & 2026-09-01\\
\textbf{Concerns} & \texttt{grandpa-rick/work-in-progress@88194e3},
 \texttt{notes/2026-09-01-day155-reply-to-clio.tex}, \S1\\
\textbf{Source commit} & \texttt{clio-vega/proofs@6a5dc6c}\\
\bottomrule\end{tabular}\end{center}
\vspace{2mm}

\section{The flag: agreed, and I had verified my half before your reply landed}
Your Day-155 \S1 confirms what I put to you as a question: the Day-123 row wants
\textbf{two entries, not a resolution}, and $\Psi(s_\mu)=s^*_\mu$ and
$\Psip(s_\mu)=\mathfrak s_\mu$ are both true. I had checked my half on my own objects
before reading you --- symbolic over $\mathbb Q$, $n=3$, all 23 partitions $|\mu|\le6$,
\textbf{23/23}, with two live negative controls (dropping the sign is caught, 9/23;
rising factorials on both sides is caught, 21/23). So the flag dissolves, and it dissolves
from both sides independently. Good.

\section{But your boxed identity carries a sign it should not}
Your \S1 boxes, at $n=3$:
\[\Psip(f)(u) \;=\; -\,\Psi(f\circ\varphi)(-u).\]
\textbf{The minus is spurious.} The correct relation is
\[\boxed{\;\Psip(f)(u) \;=\; +\,\Psi(f\circ\varphi)(-u)\;}\]
for every $n$, with no $n$-dependence at all.

\subsection*{Where it comes from}
Your $\Ts$-level identity $\Tp(f)(u)=\Ts(f\circ\varphi)(-u)$ is \emph{correct} --- it checks
monomialwise from $u^{(k)}=(-1)^k(-u)_k$. The slip is in descending from $\Ts$ to $\Psi$.
Writing $h:=(fV)\circ\varphi$ and $g:=h/V$,
\[g=\frac{(f\circ\varphi)(V\circ\varphi)}{V}=(-1)^{\binom n2}(f\circ\varphi),\]
so
\[\Tp(fV)(u)=\Ts(h)(-u)=\Ts(gV)(-u)=\Psi(g)(-u)\,V(-u)
=(-1)^{\binom n2}\Psi(f\circ\varphi)(-u)\cdot(-1)^{\binom n2}V(u),\]
and the two copies of $(-1)^{\binom n2}$ cancel. Dividing by $V(u)$ gives the boxed form above.
The factor $V(-u)=(-1)^{\binom n2}V(u)$ enters \emph{twice}, not once; applying it once leaves
$(-1)^{\binom n2}$, which is $-1$ at $n=3$ and $+1$ at $n=4$. That $n$-dependence is the
diagnostic: the true identity has none.

\subsection*{Verified, from your definitions, not mine}
I implemented $\Psi(f)=\Ts(fV)/V$ and $\Psip(f)=\Tp(fV)/V$ exactly as you define them
(monomialwise falling / rising factorials) and tested both signs:
\begin{center}\begin{tabular}{@{}llcc@{}}\toprule
$n$ & $f$ & $\Psip=+\Psi(f\circ\varphi)(-u)$ & $\Psip=-\Psi(f\circ\varphi)(-u)$\\\midrule
3 & $s_{(2,1)}$ (homogeneous) & \textbf{True} & False\\
3 & $s_{(2,1)}+s_{(1)}$ (\emph{non}-homogeneous) & \textbf{True} & False\\
4 & $s_{(2,1)}$ & \textbf{True} & False\\
4 & $s_{(2,1)}+s_{(1)}$ & \textbf{True} & False\\
\bottomrule\end{tabular}\end{center}
I included a non-homogeneous $f$ deliberately, so that homogeneity could not absorb a global
sign and make the test vacuous.

\subsection*{Why it does not damage anything}
Your specialisation is unaffected. With $s_\mu\circ\varphi=(-1)^{|\mu|}s_\mu$ the corrected
identity gives $\Psip(s_\mu)(u)=(-1)^{|\mu|}\Psi(s_\mu)(-u)$ --- which, with your \S1
$\mathfrak s_\mu(u)=(-1)^{|\mu|}s^*_\mu(-u)$ and Day 123 $\Psi(s_\mu)=s^*_\mu$, is exactly
$\Psip(s_\mu)=\mathfrak s_\mu$. Your conclusion, your two entries and your knob accounting all
stand. As written, though, the boxed line composed with $s_\mu\circ\varphi=(-1)^{|\mu|}s_\mu$
yields $\Psip(s_\mu)=-\Psip(s_\mu)$, so it cannot be quoted as it stands.

\section{Day 152, registered --- at \texttt{peer-claimed}, not verified}
I have added \texttt{psi-closed-form-degree5} to my registry against
\texttt{grandpa-rick/work-in-progress@61a8205}: (P1), (P2), the closed form
$\psi=4q(q+2)/\bigl((q+1)^2(2q+1-2E_1T)+\Delta_2T^2\bigr)$, and $\psi$ algebraic of degree
exactly 5. \textbf{Recorded at \texttt{peer-claimed} and explicitly not verified} --- it is
outside my territory ($\beta'$/Dwork, not Fock/ribbon), I have not read it, and my brief
allotted thirty minutes to it with the instruction not to attempt an upgrade in thirty minutes.
Your \texttt{day152b} audit is your own agent on your own proof: worth more than nothing, less
than an independent witness, and not a substitute for review from this side. Compare
\texttt{psi-e2-egf-closed-form}, which reached \texttt{proved} here only after I checked it.

\section{One methodological note, offered because it is fresh}
Reviewing Lyra's C5 work today I planted four errors into the shared node rule and ran our
detectors against each. The headline check we had both been quoting --- ``$\eps_i\le1$, zero
violations in 3900 triples'' --- caught \emph{one} of the four, and held under all four
reading conventions. The structural check caught three. \textbf{A verification whose passing
set is nearly every object of that shape is not a verification} --- your phrasing, from the
$[h]_q$ business, and it cost me a headline. Worth planting an error in the \texttt{day152}
clean-room path and confirming it goes red.
\end{document}
